%% design.tex
%%

%% ==============================
\section{Design Specification}
\label{sec:Design}
%% ==============================

\subsection{Reactive HPA}
Scales strictly on the currently-observed workload, with a one-step reaction delay -- the industry-standard reactive baseline NimbusGuard itself compares against.

\subsection{Aggressive PAKS (baseline reproduction)}
A feed-forward neural network (scikit-learn \texttt{MLPRegressor}, 2 hidden layers) is trained to predict the next-step workload from the previous 10 steps. At every step, the predicted load (with a small 5\% safety buffer) is converted directly into a target pod count and applied immediately -- no filtering. This reproduces the core behaviour NimbusGuard's DQN+LSTM agent exhibits: proactive, but reacting to every raw forecast, including noise.

\subsection{Stability-Aware PAKS (proposed improvement)}
Built on the exact same predictor as Aggressive PAKS, wrapped in a \texttt{StabilityAwareController} that adds two changes before a scaling action is taken:
\begin{itemize}
  \item \textbf{Exponential smoothing} of the raw predictions ($\alpha=0.4$), so a single noisy forecast cannot on its own trigger a scaling action.
  \item \textbf{Hysteresis + cooldown}: a scaling action only fires if the smoothed prediction implies a pod count that differs from the current one by more than 1 pod, \emph{and} at least 2 steps have passed since the last scaling action.
\end{itemize}
This is a direct, minimal implementation of the calmer design NimbusGuard's own future-work section gestures at (``expanding the action space,'' hybridising with more conservative scaling), applied as a simple rule rather than a second learned agent.
